The ACS records the ordered add--scale charge path needed for evaluation while
forgetting the current assignment value and some syntactic provenance.  A local
model should retain that value and should distinguish the infinitesimal additive
and multiplicative directions.  The resulting three-dimensional state space carries
both an Ehresmann connection over the two-dimensional direction plane and, in the
nondegenerate case, a contact form.

Fix real constants \(\mu\) and \(\lambda\), and let
\[
    \mathcal C=\mathbb R^3_{(u,v,a)},
    \qquad
    \pi:\mathcal C\longrightarrow\mathbb R^2_{(u,v)}.
\]
Here \(a\) is the current assignment value, while \(u\) and \(v\) parametrize
additive and multiplicative motion.  We orient \(\mathcal C\) by
\[
    du\wedge da\wedge dv.
\]

\subsection{The arithmetic contact form}

Define
\begin{equation}
    \alpha:=da-\mu\,du-\lambda a\,dv.
    \label{eq:arithmetic-contact-form}
\end{equation}
The equation \(\alpha=0\) says that a lifted base motion obeys the affine
propagation law.

\begin{proposition}[Contact nondegeneracy]
\label{prop:contact-form}
With the orientation \(du\wedge da\wedge dv\),
\begin{equation}
    \alpha\wedge d\alpha
       =\mu\lambda\,du\wedge da\wedge dv.
    \label{eq:contact-volume-sign}
\end{equation}
Hence \(\alpha\) is a contact form if and only if \(\mu\lambda\ne0\).
\end{proposition}

\begin{proof}
Since \(\mu\) and \(\lambda\) are constant,
\[
    d\alpha=-\lambda\,da\wedge dv.
\]
It follows that
\begin{align*}
    \alpha\wedge d\alpha
      &=(da-\mu\,du-\lambda a\,dv)
        \wedge(-\lambda\,da\wedge dv)\\
      &=\mu\lambda\,du\wedge da\wedge dv.
\end{align*}
This three-form is nowhere zero exactly when \(\mu\lambda\ne0\).
\end{proof}

The contact-isomorphism class is the standard local one; the arithmetic content
lies in the coordinates and in the preferred horizontal frame introduced
next.  Even in a degenerate case, \(\ker\alpha\) remains a well-defined
horizontal distribution, although it is then not contact.

\subsection{Horizontal lifts and the affine flow}

Set
\begin{equation}
    D_u:=\partial_u+\mu\partial_a,
    \qquad
    D_v:=\partial_v+\lambda a\partial_a.
    \label{eq:horizontal-arithmetic-fields}
\end{equation}
Then
\[
    \alpha(D_u)=\alpha(D_v)=0,
\]
and the two fields project respectively to \(\partial_u\) and \(\partial_v\).
Thus
\[
    \mathcal D:=\ker\alpha
       =\operatorname{span}\{D_u,D_v\}
\]
is an Ehresmann connection for \(\pi\).  Its additive and multiplicative flows
are
\begin{align}
    \Phi_u^h(u,v,a)&=(u+h,v,a+\mu h),
       \label{eq:additive-horizontal-flow}\\
    \Phi_v^k(u,v,a)&=(u,v+k,e^{\lambda k}a).
       \label{eq:multiplicative-horizontal-flow}
\end{align}

For a direction angle \(\theta\), define
\[
    D_\theta=\cos\theta\,D_u+\sin\theta\,D_v.
\]
If \((u(s),v(s),a(s))\) is tangent to \(D_\theta\), then
\[
    \frac{du}{ds}=\cos\theta,
    \qquad
    \frac{dv}{ds}=\sin\theta,
\]
and
\begin{equation}
    \frac{da}{ds}
       =\mu\cos\theta+\lambda a\sin\theta.
    \label{eq:contact-realizes-affine-flow}
\end{equation}
Thus the horizontal lift reproduces the continuous affine flow rather than
postulating a separate dynamics.

\subsection{Horizontal curvature}

\begin{theorem}[Contact curvature]
\label{thm:contact-curvature}
The preferred horizontal fields satisfy
\begin{equation}
    [D_u,D_v]=\mu\lambda\partial_a.
    \label{eq:horizontal-curvature-bracket}
\end{equation}
In particular, when \(\mu\lambda\ne0\), the horizontal distribution is not
integrable.  Its bracket defect is vertical and has constant coefficient
\(\mu\lambda\) in the arithmetic frame.
\end{theorem}

\begin{proof}
Using \eqref{eq:horizontal-arithmetic-fields},
\begin{align*}
    [D_u,D_v]
      &=[\partial_u+\mu\partial_a,
          \partial_v+\lambda a\partial_a]\\
      &=\mu\,\partial_a(\lambda a)\partial_a
        -\lambda a\,\partial_a(\mu)\partial_a\\
      &=\mu\lambda\partial_a.
\end{align*}
All remaining coordinate brackets vanish.  Since \(\partial_a\) is vertical
and is not contained in \(\ker\alpha\), the bracket is the vertical failure of
horizontal closure.
\end{proof}

The sign in \eqref{eq:horizontal-curvature-bracket} agrees with the
chronological convention of the preceding chapter.  Equivalently,
\[
    \alpha([D_u,D_v])=\mu\lambda,
    \qquad
    d\alpha(D_u,D_v)=-\mu\lambda.
\]
The two formulas have opposite signs by the identity
\(d\alpha(X,Y)=-\alpha([X,Y])\) for horizontal fields \(X,Y\).

\subsection{Open defects and closed horizontal loops}

The bracket is infinitesimal.  At finite scale there are two related, but
different, quantities.  First compare the two open lifts from \((u,v)\) to
\((u+h,v+k)\):
\begin{align*}
    \Phi_v^k\circ\Phi_u^h(u,v,a)
      &=(u+h,v+k,e^{\lambda k}(a+\mu h)),\\
    \Phi_u^h\circ\Phi_v^k(u,v,a)
      &=(u+h,v+k,e^{\lambda k}a+\mu h).
\end{align*}
Their assignment difference, with ``add then multiply'' taken first, is
\begin{equation}
    T_{\mathrm{open}}(h,k)
      =\mu h(e^{\lambda k}-1).
    \label{eq:exact-open-contact-defect}
\end{equation}
This is exactly the elementary ACS rectangle torsion with charge increments
\((\mu h,\lambda k)\).

Now perform the chronological four-step history
\[
    \mathsf A_h,\quad \mathsf M_k,\quad
    \mathsf A_{-h},\quad \mathsf M_{-k},
\]
where \(\mathsf A_h\) denotes the flow \(\Phi_u^h\) and
\(\mathsf M_k\) denotes \(\Phi_v^k\).  The base point returns to \((u,v)\),
and direct substitution gives
\begin{align*}
    a&\longmapsto a+\mu h
      \longmapsto e^{\lambda k}(a+\mu h)\\
     &\longmapsto e^{\lambda k}(a+\mu h)-\mu h
      \longmapsto a+\mu h(1-e^{-\lambda k}).
\end{align*}
The closed-loop vertical drift is therefore
\begin{equation}
    T_{\mathrm{closed}}(h,k)
       =\mu h(1-e^{-\lambda k})
       =e^{-\lambda k}T_{\mathrm{open}}(h,k).
    \label{eq:exact-closed-contact-drift}
\end{equation}
Thus the finite open comparison and the finite closed holonomy are not equal.
Here \(T_{\mathrm{closed}}\) is the source-normalized relative translation.
Converting it to the target frame gives
\(e^{\lambda k}T_{\mathrm{closed}}=T_{\mathrm{open}}\) exactly.

Figure~\ref{fig:contact-open-closed-commutator} keeps the two finite
constructions visually separate.  The square is drawn in the $(u,v)$ base;
the short vertical segments represent assignment fibers, not additional base
edges.

\begin{figure}[t]
\centering
\begin{tikzpicture}[
  x=1cm,y=1cm,
  openone/.style={-{Latex[length=2mm]},very thick,red!65!black},
  opentwo/.style={-{Latex[length=2mm]},very thick,blue!65!black},
  loop/.style={-{Latex[length=2mm]},very thick,green!45!black},
  fiber/.style={thick,black!55},
  every node/.style={font=\small}
]
  \node[font=\bfseries] at (-3.6,2.0) {open two-path comparison};
  \coordinate (o) at (-5,-0.45);
  \coordinate (ou) at (-3.45,-0.45);
  \coordinate (ouv) at (-3.45,1.1);
  \coordinate (ov) at (-5,1.1);
  \draw[openone] (o)--node[below,font=\scriptsize] {$\mathsf A_h$} (ou);
  \draw[openone] (ou)--node[right,font=\scriptsize] {$\mathsf M_k$} (ouv);
  \draw[opentwo] (o)--node[left,font=\scriptsize] {$\mathsf M_k$} (ov);
  \draw[opentwo] (ov)--node[below,font=\scriptsize] {$\mathsf A_h$} (ouv);
  \fill (o) circle (1.6pt) node[below left] {$(u,v)$};
  \fill (ouv) circle (1.6pt)
    node[above left,yshift=3pt,font=\scriptsize] {$(u+h,v+k)$};
  \draw[densely dotted,black!45] (ouv)--(-2.72,1.1);
  \draw[fiber] (-2.62,0.35)--(-2.62,1.45);
  \fill[blue!65!black] (-2.62,0.62) circle (1.7pt)
    node[right,font=\scriptsize] {$a_{\mathsf M\mathsf A}$};
  \fill[red!65!black] (-2.62,1.2) circle (1.7pt)
    node[right,font=\scriptsize] {$a_{\mathsf A\mathsf M}$};
  \draw[{Latex[length=1.6mm]}-{Latex[length=1.6mm]},black!60]
    (-2.82,0.62)--(-2.82,1.2)
    node[midway,left,font=\scriptsize] {$T_{\mathrm{open}}$};

  \node[font=\bfseries] at (2.6,2.0) {closed chronological lift};
  \coordinate (c0) at (1.2,-0.45);
  \coordinate (cu) at (2.75,-0.45);
  \coordinate (cuv) at (2.75,1.1);
  \coordinate (cv) at (1.2,1.1);
  \draw[loop] (c0)--node[below,font=\scriptsize] {$\mathsf A_h$} (cu);
  \draw[loop] (cu)--node[right,font=\scriptsize] {$\mathsf M_k$} (cuv);
  \draw[loop] (cuv)--node[above,font=\scriptsize] {$\mathsf A_{-h}$} (cv);
  \draw[loop] (cv)--node[left,font=\scriptsize] {$\mathsf M_{-k}$} (c0);
  \fill (c0) circle (1.6pt) node[below left] {$(u,v)$};
  \draw[densely dotted,black!45] (c0)--(3.48,-0.45);
  \draw[fiber] (3.58,-0.75)--(3.58,0.55);
  \fill (3.58,-0.48) circle (1.7pt)
    node[right,font=\scriptsize] {$a$};
  \fill[green!45!black] (3.58,0.25) circle (1.7pt)
    node[right,font=\scriptsize] {$a+T_{\mathrm{closed}}$};
  \draw[{Latex[length=1.6mm]}-{Latex[length=1.6mm]},black!60]
    (3.36,-0.48)--(3.36,0.25)
    node[midway,left,font=\scriptsize] {$T_{\mathrm{closed}}$};
  \node[align=center,font=\scriptsize] at (2.45,-1.5)
    {$T_{\mathrm{closed}}=e^{-\lambda k}T_{\mathrm{open}}$\\
     same infinitesimal density $\mu\lambda$};
\end{tikzpicture}
\caption{Open defect and closed holonomy in the contact model.  The two open
horizontal lifts have the same base endpoint but different assignment values.
The four-step closed base loop returns to its starting base point with a raw
vertical drift.  These finite quantities differ by the displayed scale factor;
converting the closed drift to the target frame gives the open defect, while
their infinitesimal densities agree as stated in the text.}
\label{fig:contact-open-closed-commutator}
\end{figure}

\subsection{The scalar horizontal differential}

For a smooth scalar field \(F\in C^\infty(\mathcal C)\), define
\begin{equation}
    \delta_HF
       :=dF-(\partial_aF)\alpha
       =(D_uF)\,du+(D_vF)\,dv.
    \label{eq:scalar-horizontal-differential}
\end{equation}
This is the ordinary differential restricted to horizontal directions and
written on the base coframe.  In particular,
\[
    \delta_Ha=\mu\,du+\lambda a\,dv.
\]

For scalar fields only, define the antisymmetrized second horizontal derivative
by
\begin{equation}
    \mathcal R_H(F)
      :=(D_uD_vF-D_vD_uF)\,du\wedge dv.
    \label{eq:horizontal-scalar-curvature-definition}
\end{equation}
This notation deliberately avoids presupposing an extension of \(\delta_H\) to
a graded complex.

\begin{proposition}[Curvature on scalar fields]
\label{prop:horizontal-differential-curvature}
For every \(F\in C^\infty(\mathcal C)\),
\begin{equation}
    \mathcal R_H(F)
      =\mu\lambda(\partial_aF)\,du\wedge dv.
    \label{eq:horizontal-scalar-curvature}
\end{equation}
In particular,
\[
    \mathcal R_H(a)=\mu\lambda\,du\wedge dv.
\]
\end{proposition}

\begin{proof}
By definition and Theorem~\ref{thm:contact-curvature},
\[
    \mathcal R_H(F)
       =[D_u,D_v](F)\,du\wedge dv
       =\mu\lambda(\partial_aF)\,du\wedge dv.
\]
Taking \(F=a\) proves the final formula.
\end{proof}

One may use \(\delta_H^2F\) as shorthand for
\(\mathcal R_H(F)\) after declaring this antisymmetrization convention.  It
does not mean that \(\delta_H\) has been defined on horizontal forms, nor does
it assert the existence of a nilpotent differential complex.  No such graded
calculus is required for the foundational result.

\subsection{Three levels of the same order defect}

\begin{proposition}[Finite and infinitesimal comparison]
\label{prop:torsion-curvature-synthesis}
For the elementary rectangle with additive side \(\mu h\) and logarithmic
multiplicative side \(\lambda k\), the following statements hold.
\begin{enumerate}
    \item The open affine defect equals the ACS weighted area exactly:
    \[
        T_{\mathrm{open}}(h,k)
          =\int_{[0,\mu h]\times[0,\lambda k]}
              e^{\lambda k-M}\,dA\wedge dM
          =\mu h(e^{\lambda k}-1),
    \]
    with oriented-chain interpretation when a side is negative.
    \item The closed chronological lift has exact drift
    \[
        T_{\mathrm{closed}}(h,k)
          =e^{-\lambda k}T_{\mathrm{open}}(h,k).
    \]
    Thus \(T_{\mathrm{closed}}\) is the source-normalized relative translation;
    its target-frame conversion agrees with the open defect, but the raw finite
    values generally differ.
    \item Both defects have the common infinitesimal density
    \[
        \lim_{(h,k)\to(0,0)}
          \frac{T_{\mathrm{open}}(h,k)}{hk}
        =\lim_{(h,k)\to(0,0)}
          \frac{T_{\mathrm{closed}}(h,k)}{hk}
        =\mu\lambda,
    \]
    where the quotients are taken for \(hk\ne0\).  This is the coefficient of
    \([D_u,D_v]\) and of \(\mathcal R_H(a)\).
\end{enumerate}
\end{proposition}

\begin{proof}
The first identity is Theorem~\ref{thm:torsion-stokes} applied to the elementary
rectangle, and its integral evaluates to
\[
    \mu h\int_0^{\lambda k}e^{\lambda k-M}\,dM
       =\mu h(e^{\lambda k}-1).
\]
The second identity is \eqref{eq:exact-closed-contact-drift}.  Finally,
\[
    \frac{e^{\lambda k}-1}{k}\longrightarrow\lambda,
    \qquad
    \frac{1-e^{-\lambda k}}{k}\longrightarrow\lambda,
\]
which proves the two limits.  The curvature identifications follow from
Theorem~\ref{thm:contact-curvature} and
Proposition~\ref{prop:horizontal-differential-curvature}.
\end{proof}

The synthesis has three layers, not one finite equality: exact target-frame open
affine torsion is measured by ACS weighted area; source-normalized closed horizontal
holonomy differs by an exact frame-conversion factor; and contact curvature is their common
infinitesimal limit.  The contact distribution alone supplies neither a
preferred complex structure nor a full analytic function theory.  The present
paper uses only the horizontal connection, its scalar differential, and its
curvature.
